\section*{Hoofdstuk \ref{ch:b3:remote-sensing-gis} --- Teledetectie en GIS}

\begin{solution}{exo:b3:remote-sensing-gis:1}
Veld: $(0.42 - 0.06)/(0.42 + 0.06) = 0.75$, dichte groene vegetatie, een gewas
in volle bladstand. Geploegd veld: $(0.28 - 0.21)/(0.28 + 0.21) = 0.14$, vrijwel
kale grond, een bodem die licht helderder wordt richting het infrarood.
\end{solution}

\begin{solution}{exo:b3:remote-sensing-gis:2}
(a)~$185000/30 \approx 6170$ pixels; $185000/15 \approx 12300$ in de
panchromatische band. (b)~$2^{12} = 4096$ niveaus; $2^8 = 256$. (c)~16 dagen
alleen; met Landsat~9 in dezelfde baan, acht dagen later, 8 dagen.
\end{solution}

\begin{solution}{exo:b3:remote-sensing-gis:3}
(a)~Eén interferentiestreep is $\lambda/2 = \qty{27.7}{mm}$; zes strepen zijn
\qty{166}{mm} verandering langs de \omterm{def:b3:remote-sensing-gis:insar}{zichtlijn}, hier van de satelliet af, omdat
het bekken daalt. (b)~$56.6/2 = \qty{28.3}{mm}$.
\end{solution}

\begin{solution}{exo:b3:remote-sensing-gis:4}
(a)~$t = 2d/c = 2\times\num{5e5}/\num{3.00e8} = \qty{3.34}{ms}$.
(b)~$c\,\delta t/2 = \num{3.00e8}\times\num{8e-10}/2 = \qty{0.12}{m}$; een
oppervlak dat door veel fotonen wordt bemonsterd, wordt veel nauwkeuriger
gelokaliseerd, de precisie neemt af met de wortel uit hun aantal.
\end{solution}

\begin{solution}{exo:b3:remote-sensing-gis:5}
(a)~In SI-eenheden is $L = \qty{9.0e6}{W.m^{-2}.sr^{-1}.m^{-1}}$. Met $c_1 =
2hc^2 = \qty{1.191e-16}{W.m^2.sr^{-1}}$, $c_2 = \qty{0.014388}{m.K}$ en
$\lambda^5 = (\num{1.09e-5})^5 = \num{1.539e-25}$, geldt $c_1/(\lambda^5L) = 86.0$,
en $T_b = c_2/(\lambda\ln 87.0) = \qty{295.6}{K}$. (b)~De kinetische temperatuur
volgt uit $0.98\,B_\lambda(T) = L$, dus $B_\lambda(T) = L/0.98$: dezelfde
formule geeft $T = \qty{296.9}{K}$, \qty{1.3}{K} meer.
\end{solution}

\begin{solution}{exo:b3:remote-sensing-gis:6}
(a)~$R = 693/\cos\ang{35} = \qty{846}{km}$. (b)~Reële antenne: $R\lambda/L =
\num{8.46e5}\times0.0555/12.3 = \qty{3.8}{km}$; synthetische apertuur: $L/2 =
\qty{6.2}{m}$, ongeveer 600 keer fijner. (c)~De synthetische apertuur is de
lengte van het spoor waarover een punt in de bundel blijft, $R\lambda/L =
\qty{3.8}{km}$.
\end{solution}

\begin{solution}{exo:b3:remote-sensing-gis:7}
$\partial z/\partial x = (418 - 405)/60 = 0.217$ (oost minus west),
$\partial z/\partial y = (418 - 401)/60 = 0.283$ (noord minus zuid);
$|\nabla z| = 0.357$, een helling van $\arctan 0.357 = \ang{19.6}$. Het terrein
stijgt naar het noordoosten, dus het is naar het zuidwesten gericht: $A =
\operatorname{atan2}(-0.217, -0.283) = \ang{217}$. Reliëfschaduw: $\cos\ang{45}
\cos\ang{19.6} + \sin\ang{45}\sin\ang{19.6}\cos(\ang{315} - \ang{217}) =
0.666 - 0.033 = 0.63$, iets donkerder dan vlak terrein (0.71): de helling
keert zich een weinig af van een lichtbron in het noordwesten.
\end{solution}

\begin{solution}{exo:b3:remote-sensing-gis:8}
Rijtotalen 52, 48, 38, 22; kolomtotalen 54, 46, 38, 22; diagonaal 133 van
160: totale nauwkeurigheid \qty{83.1}{\percent}. Gebruikersnauwkeurigheden: water
$48/52 = 0.92$, bos $37/48 = 0.77$, gewassen $29/38 = 0.76$, stad
$19/22 = 0.86$. Producentennauwkeurigheden: water $48/54 = 0.89$, bos $37/46 = 0.80$,
gewassen $29/38 = 0.76$, stad $19/22 = 0.86$. $p_e = (52\cdot54 + 48\cdot46 +
38\cdot38 + 22\cdot22)/160^2 = 0.271$, $\kappa = (0.831 - 0.271)/(1 -
0.271) = 0.77$. Gewassen worden het vaakst gemist: 9 van de 38 echte gewaspixels
zijn elders terechtgekomen, vooral in het bos.
\end{solution}

\begin{solution}{exo:b3:remote-sensing-gis:9}
(a)~$L = 0.04\,B_\lambda(1443) + 0.96\,B_\lambda(290)$ bij
\qty{10.9}{\micro\metre}, \qty{28.6}{W.m^{-2}.sr^{-1}.\micro\metre^{-1}},
waarvan de \omterm{def:b3:remote-sensing-gis:brightness}{helderheidstemperatuur} \qty{396}{K} is, tegenover een gemiddelde temperatuur van
$0.04\times1443 + 0.96\times290 = \qty{336}{K}$. (b)~Bij
\qty{2.2}{\micro\metre} geldt $B_\lambda(1443)/B_\lambda(290) = \num{6.8e7}$:
de lava domineert de pixel volledig, en de \omterm{def:b3:remote-sensing-gis:brightness}{helderheidstemperatuur} is
\qty{845}{K}. De kortgolvige infraroodband detecteert kleine hete plekken
veel beter; de thermische band, gevoelig voor de hele pixel, is degene die
het terrein eromheen meet.
\end{solution}

\begin{solution}{exo:b3:remote-sensing-gis:10}
Met de satelliet naar het oosten, $d_1 = u_E\sin\theta + u_z\cos\theta$; naar het
westen, $d_2 = -u_E\sin\theta + u_z\cos\theta$. Optellen en aftrekken geeft
$u_z = (d_1 + d_2)/(2\cos\theta) = 102/1.638 = \qty{62.3}{mm}$ en $u_E =
(d_1 - d_2)/(2\sin\theta) = 26/1.147 = \qty{22.7}{mm}$. Eén baan levert één
vergelijking voor twee onbekenden: elke mengvorm van opheffing en oostwaartse
beweging met dezelfde projectie op de \omterm{def:b3:remote-sensing-gis:insar}{zichtlijn} geeft dezelfde strepen.
\end{solution}

\begin{solution}{exo:b3:remote-sensing-gis:11}
(a)~Volgens \cref{thm:b3:remote-sensing-gis:fringe} vereist $\Delta\phi =
4\pi\Delta r/\lambda < \pi$ dat $\Delta r < \lambda/4 =
\qty{13.9}{mm}$. (b)~$u_z = k z_0(r^2 + z_0^2)^{-3/2}$ met $k = (1 -
\nu)\Delta V/\pi$; $\mathrm du_z/\mathrm dr = -3kz_0r(r^2 + z_0^2)^{-5/2}$,
waarvan de afgeleide nul is waar $(r^2 + z_0^2) = 5r^2$, bij $r = z_0/2$. Dan
\[
\left|\frac{\mathrm du_z}{\mathrm dr}\right| = \frac{3k}{2z_0^3}\Big(\frac54\Big)^{-5/2}
= \frac{0.859\,k}{z_0^3} = \frac{0.859\times\num{2.39e6}}{\num{6.4e10}} = \num{3.2e-5},
\]
\qty{32}{mm} per kilometer. (c)~$13.9/32 = \qty{0.43}{km}$: pixels kleiner dan
ongeveer \qty{430}{m}, wat makkelijk haalbaar is met die van Sentinel-1.
\end{solution}

\begin{solution}{exo:b3:remote-sensing-gis:12}
(a)~Voor $N, R \ge 0$ die niet beide nul zijn, geldt $|N - R| \le N + R$, dus $|\mathrm{NDVI}|
\le 1$; door $N$ en $R$ te vervangen door $aN$ en $aR$ wordt de teller en de noemer
met $a$ vermenigvuldigd. (b)~Zonder: $(0.502 - 0.051)/(0.502 + 0.051) =
0.816$; met: $(0.512 - 0.081)/(0.512 + 0.081) = 0.727$. (c)~De padradiantie verlaagt
de index van groene vegetatie, hier met ongeveer een tiende, en wel met een
bedrag dat verandert met de nevel: een kaart van ongecorrigeerde beelden zou
het weer net zoveel tonen als de vegetatie, en vergelijkingen tussen data
zouden zinloos zijn.
\end{solution}

\begin{solution}{pb:b3:remote-sensing-gis:1}
\textbf{1.} Bos $(0.38 - 0.04)/0.42 = 0.81$; lavastroom van 1890 $0.23/0.37 =
0.62$; van 1960 $0.10/0.32 = 0.31$; van 2005 $0.02/0.26 = 0.08$.

\textbf{2.} 1890, dan 1960, dan 2005: de oudste lavastroom is het groenst,
in de volgorde van de ouderdommen; ze heeft het bos nog niet bereikt.

\textbf{3.} Kaal gesteente weerkaatst in het nabije infrarood iets meer dan in het rood,
zoals het bruine zand in de figuur van het hoofdstuk; korstmossen en de eerste
planten voegen weinig toe: 0.08 is de waarde van vrijwel kale bodem.

\textbf{4.} Een helling die van het licht af is gekeerd, verzwakt beide banden in dezelfde
verhouding; de verhouding heft dit op, terwijl de reflectie in het nabije infrarood alleen
een beschaduwd bos voor schaarse vegetatie zou aanzien.

\textbf{5.} Volgens \cref{prop:b3:remote-sensing-gis:planck-inversion} met
$\lambda = \qty{10.9}{\micro\metre}$: het plateau \qty{291.1}{K}, de
fumarolepixel \qty{302.3}{K}.

\textbf{6.} $0.97\,B_\lambda(T) = \num{8.40e6}$ geeft $T = \qty{293.0}{K}$.

\textbf{7.} $B_\lambda(\qty{370}{K}) =
\qty{22.48}{W.m^{-2}.sr^{-1}.\micro\metre^{-1}}$; $f = (9.95 - 8.40)/(22.48
- 8.40) = 0.110$, ongeveer \qty{1100}{m^2} van het pixel van \qty{e4}{m^2}.

\textbf{8.} De radiantie is een convexe functie van de temperatuur
(\cref{prop:b3:remote-sensing-gis:mixed}): het hete deel levert meer radiantie op
dan zijn aandeel in de temperatuur doet verwachten, zodat het pixel heter lijkt dan zijn gemiddelde,
$0.110\times370 + 0.890\times291.1 = \qty{299.8}{K}$ tegenover
\qty{302.3}{K}.

\textbf{9.} De fase van elke echo hangt af van het exacte pad; een grote
verandering van gezichtspunt verandert de som van de echo's van de verstrooiers, en de
\omterm{def:b3:remote-sensing-gis:insar}{interferogram} verliest zijn coherentie, en er komt een fase bij door het reliëf.
$48/12 = 4$ cycli van één satelliet.

\textbf{10.} $5.4\times27.73 = \qty{149.7}{mm}$ richting de satelliet.

\textbf{11.} $u_z = 149.7/\cos\ang{35} = \qty{182.8}{mm}$.

\textbf{12.} De grond beweegt zowel naar buiten als omhoog; aan de oostkant is de
naar buiten gerichte beweging richting de satelliet en telt op bij de opheffing, aan de westkant
trekt ze ervan af (\cref{prop:b3:remote-sensing-gis:mogi-los}). De verandering
neemt sneller af naar het westen, waar de interferentiestrepen dicht op elkaar liggen; het patroon is
naar het oosten uitgerekt.

\textbf{13.} Op die lijn is de horizontale beweging noord--zuid,
loodrecht op de kijkrichting, dus de radar ziet alleen $u_z\cos\theta$, van de
vorm van de opheffing. De opheffing halveert bij $r = 0.766\,z_0$, dus $z_0 =
3.4/0.766 = \qty{4.44}{km}$.

\textbf{14.} De helft van 5.4: 2.7 interferentiestrepen.

\textbf{15.} Een vervorming is na schaling met de tijdsintervallen op beide interferogrammen dezelfde,
of ze groeit van paar tot paar door; een
atmosferisch patroon verandert van datum tot datum en volgt vaak het
reliëf. Een tweede paar dat hetzelfde concentrische patroon toont, bevestigt de
vervorming.

\textbf{16.} $\Delta V = \pi u_z(0)z_0^2/(1 - \nu) = \pi\times0.1828\times
4436^2/0.75 = \qty{1.51e7}{m^3}$.

\textbf{17.} $\num{1.51e7}/48 = \qty{3.1e5}{m^3}$ per dag, \qty{1.1e8}{m^3}
per jaar als het tempo aanhield.

\textbf{18.} $u_r \propto r/(r^2 + z_0^2)^{3/2}$ is het grootst waar haar afgeleide
verdwijnt, $r^2 + z_0^2 = 3r^2$, $r = z_0/\sqrt2 = \qty{3.14}{km}$;
daar $u_r = \qty{70}{mm}$, \qty{38}{\percent} van de centrale opheffing.

\textbf{19.} Straal \qty{1.48}{km}, volume $\frac43\pi(1.48)^3 =
\qty{13.6}{km^3}$; de verandering is $\num{1.51e7}/\num{1.36e10} =
\qty{0.11}{\percent}$ ervan.

\textbf{20.} Nee: een tiende van een interferentiestreep is \qty{2}{\percent} van het aantal, en
de diepte, afgelezen uit een halve breedte van \qty{3.4}{km}, misschien tot
\qty{0.1}{km}, gaat kwadratisch mee, \qty{6}{\percent}; de \qty{3}{\percent}
correctie valt binnen die onzekerheid.

\textbf{21.} De opwarming van het fumarolenveld (deel~II) in de tijd, een
toename van kleine aardbevingen onder de bult, een verandering in het vrijgekomen
gas, en de opheffing gemeten met satellietpositionering of waterpassing op een
paar punten van de flank.

\textbf{22.} Het reservoir onder Mount Halvren groeide met
\textbf{15~miljoen~\unit[text-series-to-math=true]{m^3}} ($\num{1.51e7}\,\unit{m^3}$).
\end{solution}
